\documentclass[conference]{IEEEtran}
\IEEEoverridecommandlockouts
\usepackage{amsmath,amssymb}
\usepackage{booktabs}
\usepackage{graphicx}
\usepackage{url}
\usepackage{microtype}
\usepackage{array}
\usepackage{enumitem}

\title{Neurofluxion: Design and Implementation of an Interactive Environment for Inspectable Neural-Network Computation}

\author{\IEEEauthorblockN{Rayyan Ahmed Khan}
\IEEEauthorblockA{Department of Computer Science and Engineering\\
HKBK College of Engineering, Bengaluru, India}}

\begin{document}
\maketitle

\begin{abstract}
Neural-network development tools typically optimize for model construction, training, and prediction, while detailed inspection of the computation that produces an output is distributed across framework APIs, profilers, debuggers, and specialized visualization systems. This paper presents Neurofluxion, an implementation-oriented full-stack environment that unifies model exploration, framework-backed execution tracing, and an independently implemented numerical simulator under a common interactive inspection workflow. The system is implemented as a React/Vite/TypeScript client and a FastAPI/Python backend. Its framework path uses TensorFlow/Keras for real model execution and exposes intermediate outputs, structural metadata, bounded tensor payloads, and selected architecture-specific internal details. Its simulator uses NumPy to construct supported sequential neural-network configurations and expose forward computation, backward computation, parameters, gradients, and training state. A model registry provides lazy loading, preprocessing metadata, checkpoint provenance, and availability status for real pretrained models. We describe the architecture, component contracts, execution semantics, provenance model, frontend workflow, and reproducibility infrastructure in implementation detail. Validation is organized as executable contracts rather than model-accuracy benchmarking: fixed finite-difference graph oracles cover Dense, Conv2D, RNN, LSTM, GRU, Embedding--Attention--RNN--Dense, and Conv2D--BatchNorm--Flatten--Dense configurations; one-batch update contracts exercise representative recurrent, convolutional/normalization, and embedding/attention/recurrent paths; adversarial browser tests preserve known asynchronous provenance defects as expected failures; and a live OpenAlex/Crossref export provides a reproducible literature-search artifact. On the frozen source revision, the live export produced 2,400 raw records and 2,074 automated deduplicated candidates with zero query failures. The resulting evidence supports a bounded implementation contribution and a reproducible software artifact, not universal numerical correctness, superior visualization, or measured human benefit.
\end{abstract}

\begin{IEEEkeywords}
neural-network visualization, interactive machine learning, neural-network simulation, execution tracing, model inspection, software architecture, reproducible software artifacts
\end{IEEEkeywords}

\section{Introduction}
Neural networks are commonly consumed through a narrow interface: an input enters a model and a prediction exits. This abstraction is appropriate for deployment, but it suppresses the intermediate computation that is often the object of debugging, experimentation, education, and systems research. A practitioner may need to know which layer produced a value, which tensor was actually returned by the framework, how a parameter contributed to a local computation, how a recurrent state evolved, or whether a displayed quantity is a complete tensor, a bounded sample, or a derived statistic. Interactive visualization systems have established the value of exposing architecture and dataflow, while modern machine-learning frameworks provide increasingly capable execution and instrumentation primitives \cite{wongsuphasawat2018,deepeyes2018,net2vis2021}.

Neurofluxion addresses this engineering problem by implementing an end-to-end environment in which executable computation and interactive inspection are designed as one workflow. The system deliberately contains two computational paths. The framework path executes selected TensorFlow/Keras models and captures intermediate outputs. The simulator path independently constructs supported networks using NumPy and exposes their forward and backward computations. A React frontend provides the interaction layer over both paths, while a FastAPI backend provides model, execution, simulation, training, device, and metadata services.

The central architectural constraint is provenance. A framework trace and a simulator trace may describe conceptually similar operations, but they are not treated as the same computation. Registry checkpoints are not implicitly imported into the simulator; simulator parameters are not implicitly represented as framework weights; and a visualization payload is not treated as a complete tensor when the implementation has bounded or sampled it. This separation is important for a software artifact whose primary contribution is inspection: the system should make it possible to distinguish what was executed from what was merely represented.

This work is therefore positioned as an implementation/system paper. It does not introduce a new learning algorithm, claim a new model architecture, or evaluate predictive accuracy against state-of-the-art models. The contribution is the design and implementation of an inspectable computational environment together with an evidence-oriented validation strategy.

The principal contributions are:
\begin{itemize}[leftmargin=*]
    \item a full-stack architecture connecting framework-backed inference, an independently implemented numerical simulator, training services, model metadata, and interactive inspection;
    \item an explicit execution-to-inspection contract covering architecture selection, execution, trace capture, stage selection, tensor inspection, mathematical views, and architecture-specific internal details;
    \item a real pretrained-model registry with lazy loading, family-specific preprocessing, availability states, caching, and checkpoint provenance;
    \item a simulator implementation that exposes supported forward/backward computation and training state rather than merely rendering a framework graph;
    \item an executable validation layer based on finite-difference graph oracles, one-batch parameter-update contracts, browser provenance adversaries, and reproducible literature-search artifacts; and
    \item a claim-discipline framework that distinguishes implementation facts, independently verified evidence, derived quantities, and unresolved limitations.
\end{itemize}

The remainder of the paper focuses on how these contributions are implemented, how their contracts are validated, and where the implementation deliberately stops making claims.

\section{Positioning and Related Systems}
Interactive neural-network visualization is an established systems area. The TensorFlow Graph Visualizer introduced graph transformation and hierarchical interaction mechanisms for making large deep-learning dataflow graphs easier to explore \cite{wongsuphasawat2018}. DeepEyes explored progressive visual analytics for neural-network design and training, emphasizing iterative inspection of layers and learned patterns \cite{deepeyes2018}. Net2Vis addressed a different but related implementation problem by defining a visual grammar for automatically generating publication-oriented CNN architecture visualizations \cite{net2vis2021}. These systems demonstrate that architecture and computation become more useful when the representation is interactive rather than static.

Neurofluxion is not intended as a replacement or direct benchmark against these systems. Its implementation emphasis is the connection between executable computation and inspectable numerical state. In particular, the simulator is an execution environment rather than a renderer of an existing graph, while the framework path is deliberately retained for real model execution. The two paths converge at the interaction layer but remain computationally distinct.

Trace-oriented interfaces are also emerging in other computational domains. Graph of Trace, for example, visualizes execution traces of scientific agents and emphasizes navigation through fine-grained execution events \cite{gao2026}. Neurofluxion adopts the same broad engineering principle---make execution state navigable---but applies it to neural-network computation with more constrained semantics. Its trace records are generated by the actual Keras execution service and by a separate simulator; they are not intended to constitute a universal execution representation.

The implementation is built on established numerical and machine-learning foundations. TensorFlow supplies the framework-backed computation layer \cite{abadi2016}; Keras provides the model-level API \cite{chollet2015}; and NumPy supplies the array-programming substrate for the independent simulator \cite{harris2020}. The novelty claim of this paper is consequently not a new numerical primitive. It is the system-level integration of these components into a provenance-aware inspection workflow and the explicit validation contracts surrounding that integration.

\section{Design Objectives and System Requirements}
The implementation was organized around six requirements.

\subsection{R1: Executable computation}
The system must expose values produced by actual execution rather than replacing computation with illustrative diagrams. Prediction Mode therefore uses real checkpoints, while the Lab executes the selected Keras model. Simulator Mode executes the numerical simulator rather than rendering a static architecture.

\subsection{R2: Inspectable intermediate state}
A user must be able to move from a model or architecture to a computation stage and inspect the state associated with that stage. The implementation consequently retains layer identity, operation metadata, shapes, parameter counts, output statistics, and selected internal values where the relevant endpoint exists.

\subsection{R3: Computational separation}
Framework execution and simulation must remain distinguishable. This requirement prevents accidental claims of shared weights or numerical equivalence and makes the source of every displayed value part of the interpretation contract.

\subsection{R4: Real model provenance}
Pretrained models must be represented by explicit metadata and availability state. Weights are downloaded lazily, cached outside source control, and accompanied by provenance information rather than silently fabricated when a checkpoint is unavailable.

\subsection{R5: Bounded serialization}
Interactive interfaces cannot always return every element of a large tensor. The implementation therefore permits bounded payloads and records sampling state. A bounded payload must remain distinguishable from a complete tensor.

\subsection{R6: Reproducible validation}
The implementation must be testable independently of visual inspection. Research workflows therefore validate selected numerical graphs, training updates, browser provenance invariants, and literature-search metadata. Passing a bounded contract must not be generalized to unsupported configurations.

\section{System Architecture}
Neurofluxion is implemented as a web application with a TypeScript/React client and a Python/FastAPI server. React 18, Vite, Tailwind CSS, and Zustand provide the frontend application layer; FastAPI and Uvicorn provide HTTP services; WebSockets support training and stream interactions; TensorFlow/Keras and NumPy provide the principal computational substrates; Three.js and Recharts provide 3-D and chart-based views; and KaTeX renders mathematical expressions.

\subsection{Component decomposition}
The backend is organized into API routes, model builders, a pretrained model registry, a simulator, training services, and general services. The frontend is divided into reusable components, design-system primitives, state stores, and page-level composition. This separation is important because the model registry and the simulator have different lifecycle requirements: registry models are external checkpoints with lazy loading, while simulator networks are constructed as local computational sessions.

\begin{table}[t]
\caption{Primary implementation components and their responsibilities}
\label{tab:components}
\centering
\footnotesize
\begin{tabular}{p{0.25\columnwidth}p{0.62\columnwidth}}
\toprule
\textbf{Component} & \textbf{Implementation responsibility}\\
\midrule
Model registry & Model metadata, adapter contracts, availability, lazy loading, caching, provenance.\\
Prediction service & Maps a registry model identifier and family-specific input to a prediction payload.\\
Execution/trace service & Executes selected Keras models and serializes intermediate outputs and metadata.\\
Simulator & Builds supported sequential networks and executes forward/backward numerical computation.\\
Training services & Provides framework-backed training and simulator training through separate execution paths.\\
Frontend stores/views & Maintains interaction state and presents architecture, tensor, math, training, and model information.\\
Research workflows & Execute numerical contracts, browser audits, integrated validation, and literature export.\\
\bottomrule
\end{tabular}
\end{table}

\subsection{Dual computational paths}
The architecture is intentionally split into a framework path and a simulator path. Figure~\ref{fig:architecture} summarizes the boundary.

\begin{figure}[t]
\centering
\small
\setlength{\tabcolsep}{3pt}
\begin{tabular}{>{\centering\arraybackslash}p{0.42\columnwidth}c>{\centering\arraybackslash}p{0.42\columnwidth}}
\toprule
\textbf{Framework path} & & \textbf{Simulator path}\\
\midrule
Registry model + family input & & User-built supported sequence\\
$\downarrow$ & & $\downarrow$\\
Keras execution & & Graph builder\\
$\downarrow$ & & $\downarrow$\\
Intermediate outputs & & NumPy forward/backward engines\\
$\downarrow$ & & $\downarrow$\\
Trace/detail records & & Step-indexed records\\
\midrule
\multicolumn{3}{c}{\textbf{Shared interaction layer:} select stage/unit $\rightarrow$ inspect tensor, math, architecture, or state}\\
\bottomrule
\end{tabular}
\caption{Implementation boundary between the framework-backed execution path and the independently implemented simulator.}
\label{fig:architecture}
\end{figure}

Both paths eventually produce data consumed by the frontend, but the provenance of a value remains tied to the path that generated it. This is a design invariant rather than merely a documentation convention.

\section{Implementation of the Framework Path}
\subsection{Model registry}
Prediction Mode uses a central registry containing model-family metadata, input requirements, preprocessing information, parameter counts, source information, and runtime availability. The repository documents real pretrained models spanning MNIST multilayer perceptrons, ImageNet CNNs including MobileNetV3-Small, MobileNetV2, ResNet50, and VGG16, and an IMDB sentiment BiLSTM. The registry also records models that are unavailable or incompatible rather than presenting them as usable checkpoints.

A model adapter follows a common contract for description, loading, prediction, and availability checking. The registry resolves a model identifier to an adapter, while the prediction service remains independent of model-specific implementation details. This design makes adding a model an adapter-and-registration operation rather than a modification to the frontend's prediction logic.

Weights are downloaded lazily at first load or prediction and cached outside the repository. A small provenance record accompanies cached weights. Keras ImageNet checkpoints use the Keras cache, while registry-specific weights use the configured model cache. The repository therefore avoids committing large binary checkpoints while preserving enough metadata to identify their source.

\subsection{Family-specific input contracts}
The frontend does not send a single universal input representation. The ANN family uses a 28$\times$28 digit canvas, CNN models accept images prepared for 224$\times$224 inference, and the RNN family accepts text that is tokenized and padded according to the packaged model's requirements. The adapter owns model-specific preprocessing so that the frontend remains model-agnostic.

\subsection{Execution tracing}
The Lab execution path creates an activation model over the non-input outputs of the selected Keras model and evaluates it on a prepared input. The returned stage records include layer identity and type, operation metadata, activation information, shapes, parameter counts, timing fields, and output statistics. Convolutional and pooling results may be rearranged for visualization, and large output payloads may be bounded before JSON serialization.

The trace therefore has two distinct products: structural metadata that describes the execution stage, and numerical payloads that may be complete, bounded, or derived. The frontend is expected to preserve this distinction.

\subsection{Timing semantics}
The reported stage timing is an attribution measure rather than a kernel-level profiler measurement. If $T_k$ is the measured execution time for a cumulative model prefix ending at stage $k$, the implementation conceptually derives
\begin{equation}
\widehat{t}_k=T_k-T_{k-1}.
\end{equation}
This difference contains framework scheduling, repeated prediction overhead, and measurement noise. It is therefore appropriate for relative execution attribution inside the interactive trace but should not be interpreted as an isolated hardware-kernel duration.

\subsection{Architecture-specific detail endpoints}
The execution service also contains detail paths for selected computation types. Dense-neuron inspection can expose incoming values, weights, bias, pre-activation, and activation. Convolution inspection uses spatial and channel indices. LSTM inspection exposes timestep- and gate/state-related values. These are intentionally specialized endpoints rather than a claim that every layer has a universal symbolic explanation.

\section{Implementation of the Numerical Simulator}
The simulator is implemented as a separate numerical environment using NumPy. It contains layer implementations, a graph builder, forward and backward engines, optimization/training components, and session-related state.

\subsection{Graph construction}
The audited builder constructs a validated sequential composition of supported layer objects and propagates shapes through the configured sequence. Observed dispatch includes dense/output layers, batch normalization, Conv2D, max/average pooling, flatten, vanilla RNN, LSTM, GRU, embedding, attention, and residual-related operations.

The builder maintains layer instances, parameter-bearing layer mappings, activation/pre-activation collections, parameter and FLOP totals, and an architecture identifier. The architecture identifier is derived from a sequence of layer neuron counts and is not a complete serialization of every layer configuration. Accordingly, the system is described as a supported sequential construction path rather than as a general arbitrary-edge graph engine.

\subsection{Forward execution}
For a configured sequence, the simulator implements the abstract relation
\begin{equation}
h^{(0)}=x,\qquad h^{(\ell)}=f_{\ell}(h^{(\ell-1)};\theta_{\ell}).
\end{equation}
For dense layers, the trace can expose a finer decomposition,
\begin{equation}
z_j=\sum_i w_{ji}h_i+b_j,\qquad h_j=\phi(z_j),
\end{equation}
including matrix multiplication, bias addition, and activation records. Mixed or non-dense paths instead expose layer-level records. The difference is an implementation property and means that trace-step counts are not a universal measure of computational depth.

The forward-step endpoint regenerates the full forward sequence and then selects the requested step. It therefore implements indexed replay/inspection rather than incremental execution. This behavior is relevant to both the user experience and the interpretation of timing or computational-cost claims.

\subsection{Backward execution and training}
The simulator exposes reverse traversal through the configured layers. In the dense-specific path, the implementation computes a loss, derives layer deltas, and forms weight and bias gradients. The generic path initializes an output gradient and invokes each instantiated layer's backward method in reverse order. These paths have different contracts and are not collapsed into a single universal backpropagation implementation claim.

The simulator also exposes parameter-update behavior through its training components. The research suite exercises this behavior on representative recurrent, convolutional/normalization, and embedding/attention/recurrent configurations. These contracts test the expected update on fixed batches rather than convergence or generalization.

A notable numerical boundary is the loss/activation relationship. The audited dense path handles MSE differently from other configured loss names when forming the final delta. For this reason, numerical claims in this paper are tied to explicitly tested configurations and are not generalized to arbitrary activation/loss combinations.

\section{Frontend and Interaction Implementation}
The frontend is designed around the lifecycle
\begin{equation*}
\textit{select/construct}\rightarrow\textit{execute}\rightarrow\textit{capture}\rightarrow\textit{select}\rightarrow\textit{inspect}.
\end{equation*}
Prediction Mode selects a real registry model and family-specific input. Lab Mode executes a framework model and exposes its trace. Simulator Mode constructs a custom network, executes it, and exposes forward/backward state. Training surfaces consume a separate training channel and present live telemetry.

The principal inspection surfaces include architecture views, tensor views, mathematical expressions, model metadata, weight inspection, training telemetry, and selected internal-detail panels. Three.js supports topology-oriented 3-D views, Recharts supports numerical charts, and KaTeX renders layer equations.

The frontend is not treated as an explainability algorithm. It is an inspection client over executable state. This distinction is important because some values are directly returned by computation, some are bounded or sampled for serialization, and some are derived for visualization. The user-facing representation must therefore preserve the semantics supplied by the backend.

\subsection{Training and stream separation}
The implementation contains multiple communication channels with deliberately different semantics. The \texttt{/train} WebSocket delegates to the framework training manager. Simulator training uses a separate simulator session/training path. The \texttt{/stream} channel is a synthetic topology-and-metrics demonstration with hard-coded identifiers, perturbations, and formula-derived curves; it is not live model execution. These channels are not presented as one shared training state.

\subsection{Asynchronous provenance}
The research suite contains browser tests for stale inspection state after graph replacement and out-of-order responses. These tests currently reproduce known runtime defects and remain expected failures. They are retained as executable specifications so that a future runtime fix can be evaluated against the same invariants. No runtime fix is claimed in this manuscript because application-source modification is outside the current research-only scope.

\section{Validation and Evidence Methodology}
The validation strategy treats the software artifact as the research object. Instead of relying only on screenshots or end-to-end demonstrations, the repository contains independent contracts that exercise numerical and interaction boundaries.

\subsection{Finite-difference graph oracles}
For selected small graphs, analytical gradients are compared with central finite differences. The merged research suite covers fixed configurations containing Dense, Conv2D, vanilla RNN, LSTM, GRU, Embedding--single-head Attention--RNN--Dense, and Conv2D--BatchNorm--Flatten--Dense compositions.

For a scalar objective $L(\theta)$ and parameter element $\theta_i$, the finite-difference reference is
\begin{equation}
\frac{\partial L}{\partial\theta_i}\approx\frac{L(\theta_i+\epsilon)-L(\theta_i-\epsilon)}{2\epsilon}.
\end{equation}
Each oracle fixes its graph, input, objective, data type, seed, perturbation, and tolerance. A pass therefore establishes agreement for the declared configuration rather than arbitrary operator correctness.

\subsection{One-batch training contracts}
The training contracts complement gradient checks by exercising the update path itself. They cover LSTM/GRU recurrent configurations, Conv2D/BatchNorm configurations, and Embedding/Attention/RNN/Dense configurations. The purpose is to verify the expected parameter transition on a fixed batch. These experiments do not establish convergence, generalization, or optimizer superiority.

\subsection{Browser provenance contracts}
The browser suite contains adversarial cases in which the graph changes while an inspection request is pending and cases in which responses arrive out of order. The current expected-failure status is evidence of reproducibility of the defects, not evidence that the defects are fixed. This distinction is preserved in the paper and in the research ledger.

\begin{table}[t]
\caption{Research evidence and claim boundaries}
\label{tab:evidence}
\centering
\footnotesize
\begin{tabular}{p{0.28\columnwidth}p{0.18\columnwidth}p{0.40\columnwidth}}
\toprule
\textbf{Evidence} & \textbf{Status} & \textbf{Supported interpretation}\\
\midrule
Graph finite differences & Passed & Fixed graph/input/objective agreement only.\\
Recurrent update contracts & Passed & Fixed one-batch update behavior.\\
Conv2D/BatchNorm update contract & Passed & Fixed architecture and batch.\\
Embedding/Attention/RNN/Dense update contract & Passed & Fixed composition and update path.\\
Browser provenance adversaries & Expected failure & Known stale-state/race behavior remains reproducible.\\
Literature metadata export & Passed & Reproducible export; screening is not complete.\\
\bottomrule
\end{tabular}
\end{table}

\section{Reproducibility and Software-Artifact Record}
IEEE guidance emphasizes that a reproducible article should provide enough methodological detail for an independent researcher to follow the same procedure, and encourages sharing code and data \cite{ieeerepro}. Neurofluxion adopts this principle at the artifact level. The repository contains source code, test suites, dedicated research workflows, validation ledgers, and a literature-search exporter. The implementation is also structured so that model checkpoints are external artifacts rather than opaque files embedded in the source tree.

The evidence baseline for the current manuscript is main revision \texttt{50de1fb7328d82c8f4799bbf3feffff8fd713216}. The live OpenAlex/Crossref export was produced from this revision in GitHub Actions using CPython 3.11.17 on Ubuntu 24.04.5. Twelve queries completed successfully, producing 2,400 raw records and 2,074 automated deduplicated candidates, with zero query failures and zero partial-query failures.

These counts describe a reproducible metadata export, not a completed systematic review. Human screening, full-text extraction, duplicate adjudication beyond automated deduplication, and citation snowballing remain separate activities. The paper therefore does not use the export counts to make claims of comprehensive literature coverage or novelty priority.

For computational reproduction, framework demonstrations require the selected checkpoint, preprocessing contract, environment, device, and input. Simulator demonstrations require the graph specification, parameters, input, objective, seed, data type, and execution path. The repository's test workflows additionally record exact source revisions and validation outcomes so that a result can be associated with a concrete code state.

\section{Implementation Tradeoffs}
The architecture intentionally accepts several tradeoffs.

\subsection{Independent simulator versus direct framework reuse}
Reusing TensorFlow execution everywhere would reduce implementation duplication but would also remove the independent numerical execution environment that makes small computations directly inspectable. Maintaining a separate NumPy simulator increases implementation and validation cost, but it permits explicit parameter, activation, gradient, and step inspection for supported configurations.

\subsection{Bounded tensors versus unrestricted serialization}
Returning every tensor element would improve completeness for some experiments but would make large models impractical for an interactive browser. The implementation therefore bounds selected payloads and retains sampling metadata. This favors responsiveness while preserving the distinction between a visualization payload and a complete tensor.

\subsection{Replay-based step selection versus incremental execution}
Recomputing a complete forward sequence before selecting a step simplifies implementation and keeps step records consistent with one execution path. The tradeoff is that step selection is not computationally equivalent to incremental execution and should not be interpreted as such.

\subsection{Adapter-based model registry}
A common adapter contract isolates model-specific loading and preprocessing from the frontend and prediction service. The tradeoff is that each new checkpoint family requires an explicit adapter and availability policy. This cost is preferable to silently assuming that heterogeneous checkpoints share the same input or runtime contract.

\subsection{Expected-failure tests}
Leaving known provenance defects as expected failures can appear less polished than making the continuous-integration dashboard entirely green. The implementation instead treats these tests as executable documentation of unresolved behavior. This makes the defect visible and prevents a false claim that asynchronous provenance is already correct.

\section{Discussion}
Neurofluxion's principal contribution is an implementation pattern for provenance-aware inspection. The system does not attempt to solve interpretability in the broad sense. It instead makes a narrower claim: an interactive neural-network environment can be designed so that architecture, execution state, numerical values, and their provenance remain connected throughout the user workflow.

The dual-path architecture is particularly important. Framework-backed execution provides realistic behavior of selected pretrained or Keras-defined models. The simulator provides a smaller and more explicit numerical environment in which supported computations can be inspected and independently validated. Treating these paths as equivalent would be simpler rhetorically but technically incorrect. The implementation therefore uses separation as a feature of the system contract.

The validation results should also be interpreted at the same granularity. A finite-difference oracle is strong evidence for the exact graph and objective that it tests, but it is not evidence for every possible graph. A one-batch update contract demonstrates a parameter transition, but it is not a convergence study. An expected-failure browser test demonstrates that a defect can be reproduced, but it is not a usability result. This alignment between evidence granularity and claim granularity is a central methodological property of the implementation paper.

Relative to visualization-only systems, Neurofluxion emphasizes the path from executable state to inspection. Relative to general-purpose machine-learning frameworks, it emphasizes interactive access to intermediate values and small numerical execution traces. Relative to educational simulators, it additionally integrates real framework-backed models and a model registry. These are positioning statements about implementation emphasis, not claims of superiority.

\section{Limitations and Threats to Validity}
The simulator's primary construction path is a validated sequential composition rather than an arbitrary-edge directed acyclic graph. Some layer configurations and generic backward paths therefore require additional evidence before broad correctness claims can be made.

The numerical validation suite is selective. It covers representative fixed graphs and training paths but does not exhaust every tensor shape, activation, loss, optimizer, sequence length, or layer composition. The loss/activation relationship in the dense-specific backward implementation is one example of why configuration-specific validation remains necessary.

The framework trace is not a hardware profiler. Stage timing is derived from cumulative prefix execution and includes framework overhead and measurement noise. Large tensor outputs may be bounded or sampled for serialization. A displayed value can therefore be useful for inspection without being a complete numerical dump.

The model registry depends on external checkpoint sources and local caches. Availability can change independently of the source repository. Reproducibility consequently requires checkpoint provenance and preprocessing information in addition to source code.

The current browser provenance tests reproduce stale-state and out-of-order response defects. The paper does not claim a runtime fix. A future implementation revision should turn these expected failures into passing invariants only after the runtime behavior is explicitly corrected and revalidated.

The literature stage is a reproducible metadata export rather than a completed systematic review. The current evidence is insufficient for comprehensive-coverage, novelty-priority, or superiority claims based on search counts alone.

Finally, the work does not contain a controlled human-subject study. It therefore cannot establish improvements in usability, learning outcomes, debugging speed, or educational effectiveness.

\section{Conclusion}
This paper presented the design and implementation of Neurofluxion, a full-stack interactive environment for inspecting neural-network computation. The implementation combines a TensorFlow/Keras framework path for real model execution and trace capture with an independently implemented NumPy simulator for supported numerical construction, forward/backward computation, and parameter inspection. A model registry supplies real pretrained checkpoints with explicit loading and provenance behavior, while the frontend connects architecture, tensor, mathematical, model, and training views into a common inspection lifecycle.

The research artifact is supported by bounded executable evidence: fixed finite-difference graph oracles, one-batch training-update contracts, browser provenance adversaries, and a reproducible literature-search export. The evidence is intentionally reported at the level at which it was obtained. This prevents a successful test of one graph from becoming a claim about arbitrary operators and prevents a visualization from being mistaken for a general interpretability method.

The resulting contribution is best understood as an implementation and systems artifact: a concrete architecture, execution model, inspection workflow, and validation methodology for making neural-network computation navigable while preserving provenance. Future work should address the remaining asynchronous provenance defects, broaden graph-level numerical coverage, improve complete checkpoint reproducibility, and conduct controlled user evaluation before making claims about usability or educational benefit.

\begin{thebibliography}{00}

\bibitem{abadi2016}
M. Abadi, P. Barham, J. Chen, Z. Chen, A. Davis, J. Dean, M. Devin, S. Ghemawat, G. Irving, M. Isard, M. Kudlur, J. Levenberg, R. Monga, R. Moore, D. G. Murray, S. G. Steiner, J. Tucker, V. Vasudevan, P. Warden, M. Wicke, Y. Yu, and X. Zheng, ``TensorFlow: A System for Large-Scale Machine Learning,'' in \emph{Proc. 12th USENIX Symp. Operating Systems Design and Implementation (OSDI)}, 2016, pp. 265--283.

\bibitem{chollet2015}
F. Chollet et al., ``Keras,'' 2015. [Online]. Available: \url{https://keras.io}

\bibitem{harris2020}
C. R. Harris et al., ``Array programming with NumPy,'' \emph{Nature}, vol. 585, no. 7825, pp. 357--362, 2020, doi: 10.1038/s41586-020-2649-2.

\bibitem{wongsuphasawat2018}
K. Wongsuphasawat, D. Smilkov, J. Wexler, J. Wilson, D. Man\'e, D. Fritz, D. Krishnan, G. Vi\'egas, and M. Wattenberg, ``Visualizing Dataflow Graphs of Deep Learning Models in TensorFlow,'' \emph{IEEE Trans. Visualization and Computer Graphics}, vol. 24, no. 1, pp. 1--12, 2018, doi: 10.1109/TVCG.2017.2744878.

\bibitem{deepeyes2018}
N. Pezzotti, T. H\"ollt, J. van Gemert, B. P. F. Lelieveldt, E. Eisemann, and A. Vilanova, ``DeepEyes: Progressive Visual Analytics for Designing Deep Neural Networks,'' \emph{IEEE Trans. Visualization and Computer Graphics}, vol. 24, no. 1, pp. 98--108, 2018, doi: 10.1109/TVCG.2017.2744358.

\bibitem{net2vis2021}
A. B\"auerle, C. van Onzenoodt, and T. Ropinski, ``Net2Vis: A Visual Grammar for Automatically Generating Publication-Tailored CNN Architecture Visualizations,'' \emph{IEEE Trans. Visualization and Computer Graphics}, vol. 27, no. 6, pp. 2980--2991, 2021, doi: 10.1109/TVCG.2021.3057483.

\bibitem{gao2026}
T. Gao, H. Li, J. H. Li, T. Zhao, R. Shi, W. Wang, Z. Wu, and L. Mi, ``Graph of Trace: Visualizing Execution Traces of Scientific Agents,'' in \emph{Proc. 64th Annual Meeting of the Association for Computational Linguistics, Volume 3: System Demonstrations}, 2026, pp. 297--306, doi: 10.18653/v1/2026.acl-demo.29.

\bibitem{ieeerepro}
IEEE Author Center, ``Research Reproducibility,'' IEEE, 2026. [Online]. Available: \url{https://journals.ieeeauthorcenter.ieee.org/create-your-ieee-journal-article/research-reproducibility/}

\end{thebibliography}

\end{document}
